\documentclass[11pt,letterpaper]{article}

\usepackage[spanish,es-noshorthands,es-tabla]{babel}
\usepackage{fontspec}
\setmainfont{Source Serif Pro}
\usepackage{microtype}
\usepackage{geometry}
\usepackage{enumitem}
\usepackage{xcolor}
\usepackage{csquotes}
\usepackage{xurl}
\usepackage[hidelinks]{hyperref}
\usepackage[backend=biber,style=apa,sorting=nyt,maxcitenames=2]{biblatex}
\DeclareLanguageMapping{spanish}{spanish-apa}

\geometry{margin=2.3cm}
\hbadness=10000
\setlength{\emergencystretch}{1em}
\setlength{\parindent}{0pt}
\setlength{\parskip}{0.5em}
\setlist{nosep}
\addbibresource{referencias.bib}
\AtBeginBibliography{\sloppy\small}
\setlength{\bibitemsep}{0.3em}

\definecolor{cEje}{RGB}{26,58,102}
\definecolor{cGris}{RGB}{120,116,110}
\definecolor{cFondo}{RGB}{244,242,238}

% Enunciado de cada paso, en gris y con sangría.
\newcommand{\pregunta}[1]{%
  \begin{quote}\color{cGris}\small\itshape #1\end{quote}}

% Recuadro para el prompt original.
\newenvironment{prompt}{%
  \begin{center}\begin{minipage}{0.92\textwidth}\color{black!80}\small
  \setlength{\parskip}{0.4em}}{%
  \end{minipage}\end{center}}

\begin{document}

% --- Portada -----------------------------------------------------------------
\begin{titlepage}
\centering
\vspace*{3.2cm}
{\large\scshape Inteligencia Artificial Generativa\\
Aplicada a la Transformación Digital\par}
\vspace{0.5cm}
{\normalsize Módulo 2. Actividad 2.1\par}
\vspace{2.6cm}
{\color{cEje}\rule{0.55\textwidth}{0.8pt}\par}
\vspace{0.7cm}
{\LARGE\bfseries ¿Publicaría lo que escribió la IA?\par}
\vspace{0.5cm}
{\large Las cuatro preguntas aplicadas a un artículo\\
sobre gemelos digitales para fondos de inversión\par}
\vspace{0.7cm}
{\color{cEje}\rule{0.55\textwidth}{0.8pt}\par}
\vfill
{\large Carlos Luis Rojas Aragonés\par}
\vspace{0.3cm}
{\normalsize Chief Technology Officer Program\par}
\vspace{0.3cm}
{\normalsize 7 de noviembre de 2026\par}
\vspace*{1.5cm}
\end{titlepage}

% --- Paso 1 ------------------------------------------------------------------
\section*{1. El resultado que evalué}

\pregunta{Utilice una herramienta de IA generativa de su elección para crear
un resultado relevante para su contexto profesional.}

Le pedí a un agente de Claude Code que redactara para mi blog un artículo
sobre gemelos digitales aplicados a fondos de inversión privados, a partir de
lo que estamos probando en Subaxa. El resultado fue un artículo en inglés con
cinco figuras: un diagrama de Sankey con el destino del dinero de la venta de
una propiedad y un modelo de ejemplo con 2000 simulaciones
\parencite{rojas2026twins}. Este fue el prompt:

\begin{prompt}
\itshape
Let's write an post in my blog about \enquote{Digital twins for private
investment funds}.
what I want is to briefly explain the Digital Twin concept
and the explain how a digital twin works in real life. We are experimenting
with them for our private investment funds at Subaxa and it help us
visualize different scenarios and analyze different paths we could take and
the eventual results in a safe and simulated environment. we can't disclosure
subaxa comercial secrets but we can describe it at a high level. make sure
you include images and/or charts to ilustrate and make the article easier
to read.

I want to include an extra image to illustrate a cash distribution in a
property investment. a Sankey diagram would make a lot of sense.
\end{prompt}

Lo evalúo comparando la primera versión que generó la IA con la que terminé
publicando. Las dos quedaron en el historial del repositorio del blog.

% --- Paso 2 ------------------------------------------------------------------
\section*{2. Las cuatro preguntas}

\pregunta{¿La información es precisa? ¿Está actualizada? ¿Encaja con este
contexto? ¿Qué habría que verificar antes de actuar?}

\textbf{¿Es precisa?} En su mayoría, sí. Las cifras del Sankey cuadran y las
de las simulaciones salen de un código que se puede volver a correr
\parencite{rojas2026codigo}. Pero la primera versión decía que los
inversionistas casi nunca pueden probar una estrategia antes de ejecutarla,
y eso es falso: el análisis de escenarios y las simulaciones de Monte Carlo
se usan en finanzas desde hace décadas. Queda además un detalle menor: la
figura dice 3\,\% de escenarios con pérdidas y el texto, 2,7\,\%.

\textbf{¿Está actualizada?} El concepto, sí. La historia del término es otra
cosa: la IA contó la versión clásica (Grieves en 2002, NASA en 2010) sin
mencionar que hay revisiones que la cuestionan.

\textbf{¿Encaja con el contexto?} En tono y alcance, sí. No revela nada de
Subaxa y marca cada cifra como ilustrativa. Lo que no encajaba era llamar
gemelo digital a una simulación que no está conectada a ningún fondo real,
cuando esa conexión es justamente lo que define a un gemelo
\parencite{nasem2023}.

\textbf{¿Qué verificaría antes de actuar?} Las fuentes, porque la primera
versión no tenía ninguna; que los rendimientos digan que son antes de
comisiones e impuestos; el supuesto de que las propiedades son
independientes; y que las frases sobre cómo usamos el gemelo en Subaxa
reflejen lo que de verdad hacemos.

% --- Paso 3 ------------------------------------------------------------------
\section*{3. Mi valoración como líder}

\pregunta{¿Utilizaría este resultado en un contexto empresarial real? ¿Qué
revisiones o salvaguardas aplicaría? ¿Qué riesgos surgirían si se usara sin
revisión?}

Sí lo usé: el artículo está publicado. Pero no la primera versión, sino la
revisada. Para un blog, el riesgo es de reputación. Para decidir una
inversión no usaría el modelo de ejemplo, porque no está calibrado con datos
reales.

Las salvaguardas que aplicaría son leer cada afirmación como si la hubiera
escrito un analista nuevo, exigir una fuente para cada dato, publicar el
código para que las figuras se puedan reproducir y revisar que nada revele
información comercial. Si el modelo se usara para decidir, lo calibraría con
datos del fondo y le agregaría choques que afectan a todas las propiedades
a la vez, como una recesión o un cambio en los juzgados.

Sin revisión, los riesgos son publicar errores con mi nombre, que un
inversionista lea el 1,32 como una promesa de rendimiento y que alguien
concluya que diversificar elimina el riesgo de perder dinero, cuando eso solo
pasa en un modelo que ignora los choques compartidos.

% --- Paso 4 ------------------------------------------------------------------
\section*{4. Reflexión sobre el uso responsable de la IA}

\pregunta{¿Qué revela este ejercicio sobre el papel del criterio humano al
utilizar IA? ¿Cómo influirá esto en la forma en que usted o su organización
utilicen la IA en adelante?}

La IA escribió en minutos un artículo que me habría tomado mucho más tiempo,
y lo escribió bien. Justamente por eso hay que tener cuidado. Los errores no
estaban en la redacción, sino en afirmaciones que suenan razonables y que
solo se detectan si uno conoce el negocio y el tema. El criterio humano ya
no está tanto en escribir como en saber qué preguntar y qué no dar por
cierto. Un texto bien escrito genera confianza, y esa confianza es la que
hay que poner a prueba. Revisar no es desconfiar de la herramienta, es hacer
el trabajo que le toca a uno.

En adelante voy a tratar lo que produce la IA como el borrador de un analista
junior: rápido y útil, pero firmado por mí. En Subaxa quiero que cualquier
texto o modelo hecho con IA que llegue a inversionistas pase por una revisión
de fuentes y de cifras, y que los modelos que apoyen decisiones se validen
con datos reales antes de usarlos. La IA acelera el trabajo; la
responsabilidad no se delega.

\printbibliography[title={Referencias}]

\end{document}
